% =============================================================================
% CompetitiveExam/CS401/08-io-techniques-answers.tex
% Answer key for 08-io-techniques-questions.tex.
%
% Q1-Q7 are real GATE past-year questions sourced from the local GATEOverflow
% Volume 3 PYQ book (CompetitiveExam/Books/index.md, Tier-1 availability), which
% authenticates each question (year + question number printed on the page).
% The book's own answer keys are collapsed and not extractable, and web
% cross-verification was attempted but all search engines were bot-blocked in
% the build session (2026-08-11). The answers below are therefore derived and
% justified step by step from the course material (Slides/Notes, which
% page-verify against Hamacher 6th ed. Ch. 3 / Ch. 7), standing in for the
% book's collapsed key --- the same policy Lecture 07's key used.
% =============================================================================

\documentclass[11pt]{article}
\usepackage[margin=1in]{geometry}
% =============================================================================
% Preambles/header.tex — shared LaTeX/Beamer preamble for this project
%
% Use from a Beamer lecture by prepending:
%     % =============================================================================
% Preambles/header.tex — shared LaTeX/Beamer preamble for this project
%
% Use from a Beamer lecture by prepending:
%     % =============================================================================
% Preambles/header.tex — shared LaTeX/Beamer preamble for this project
%
% Use from a Beamer lecture by prepending:
%     \input{header}
% and compiling with TEXINPUTS=../Preambles:$TEXINPUTS (the /compile-latex
% skill does this automatically).
%
% PALETTE CONTRACT. Color names defined here MUST match the names in
% Quarto/theme-template.scss so Beamer and Quarto renderings use the same
% palette. The scripts/check-palette-sync.sh script enforces this.
%
% To customize: edit the HEX values below AND the matching $variable in
% Quarto/theme-template.scss, then run scripts/check-palette-sync.sh.
%
% TITLE PAGE. A formal, serif (Libertine) title-page template is defined in
% the Beamer-only block below (\setbeamertemplate{title page}) -- title,
% subtitle, a thin gold rule, then \author/\institute/\date. Frametitle and
% body text remain on Beamer's sans-serif default. No SCSS counterpart is
% needed for this (font/layout only, no new colors).
% =============================================================================

% -----------------------------------------------------------------------------
% Engine detection + encoding
%
% Our /compile-latex skill uses XeLaTeX, where inputenc is unnecessary and
% can produce encoding warnings. Only load inputenc under pdfTeX.
% -----------------------------------------------------------------------------
\usepackage{iftex}
\ifPDFTeX
  \usepackage[utf8]{inputenc}
\fi

\usepackage{amsmath, amssymb, amsthm}
\usepackage{booktabs}
\usepackage{graphicx}

% -----------------------------------------------------------------------------
% Serif face for the title page (formal/academic look). Linux Libertine is
% XeLaTeX- and pdfTeX-safe and redefines \rmdefault, so \rmfamily picks it up
% everywhere \rmfamily is used explicitly. Body text and frametitles stay on
% Beamer's own sans-serif default (unaffected) -- only the title-page
% \setbeamerfont{...} overrides below opt in to \rmfamily.
% -----------------------------------------------------------------------------
\usepackage{libertine}

% -----------------------------------------------------------------------------
% PALETTE — mirrors Quarto/theme-template.scss variable names.
% Load xcolor and define colors BEFORE anything that references them
% (notably hyperref, hypersetup, and the Beamer color assignments).
% -----------------------------------------------------------------------------
\usepackage{xcolor}

\definecolor{primary-blue}{HTML}{012169}      % headings, accents
\definecolor{primary-gold}{HTML}{B9975B}      % emphasis, borders
\definecolor{highlight-yellow}{HTML}{F2A900}  % markers, alerts
\definecolor{light-bg}{HTML}{E8EDF5}          % title / section backgrounds
\definecolor{jet}{HTML}{1A1A1A}               % body text

% Semantic colors (used in diagrams and annotations)
\definecolor{positive}{HTML}{15803D}          % good / observed / identified
\definecolor{negative}{HTML}{B91C1C}          % bad / problematic / bias
\definecolor{neutral}{HTML}{525252}           % reference / context

% Highlight accents (match .hi-* classes in the SCSS)
\definecolor{hi-slate}{HTML}{314F4F}
\definecolor{hi-green}{HTML}{2E8B57}
\definecolor{hi-red}{HTML}{C41E3A}

% Semantic aliases so skill templates and snippets can write
% \textcolor{accent}{...} without hard-coding "primary-blue".
\colorlet{accent}{primary-blue}
\colorlet{accent2}{primary-gold}

% -----------------------------------------------------------------------------
% Hyperref — loaded AFTER xcolor + palette so link colors resolve.
% -----------------------------------------------------------------------------
\usepackage{hyperref}
\hypersetup{colorlinks=true, linkcolor=primary-blue, urlcolor=primary-blue,
            citecolor=primary-blue, pdfborder={0 0 0}}

% -----------------------------------------------------------------------------
% Course code — set once per deck (after \input{header}, before \begin{document})
% via \coursecode{CS401}. Shown in the footer on every slide so a deck's course
% is identifiable without opening the title page. Empty by default.
% -----------------------------------------------------------------------------
\makeatletter
\newcommand{\coursecode}[1]{\gdef\@coursecodetext{#1}}
\gdef\@coursecodetext{}
\makeatother

% -----------------------------------------------------------------------------
% Beamer theme assignments (only apply under Beamer).
%
% We detect Beamer via \@ifclassloaded{beamer}, which is the documented,
% non-internal way. This must be wrapped in \makeatletter/\makeatother
% because \@ifclassloaded contains an @-catcode character.
% -----------------------------------------------------------------------------
\makeatletter
\@ifclassloaded{beamer}{%
  \usefonttheme{professionalfonts}
  \setbeamercolor{structure}{fg=primary-blue}
  \setbeamercolor{frametitle}{fg=primary-blue}
  \setbeamercolor{title}{fg=primary-blue}
  \setbeamercolor{subtitle}{fg=primary-gold}
  \setbeamercolor{author}{fg=primary-blue}
  \setbeamercolor{institute}{fg=neutral}
  \setbeamercolor{date}{fg=primary-gold}
  \setbeamercolor{itemize item}{fg=primary-blue}
  \setbeamercolor{itemize subitem}{fg=primary-gold}
  \setbeamercolor{alerted text}{fg=primary-gold}
  \setbeamercolor{block title}{fg=white, bg=primary-blue}
  \setbeamercolor{block body}{bg=light-bg}
  % Semantic block variants: block=definition/concept (blue), exampleblock=
  % worked example (green/positive), alertblock=key takeaway or caution (gold).
  % Keep to <=2 colored boxes per slide across ALL three variants (INV-7).
  \setbeamercolor{block title example}{fg=white, bg=positive}
  \setbeamercolor{block body example}{bg=light-bg}
  \setbeamercolor{block title alerted}{fg=white, bg=primary-gold}
  \setbeamercolor{block body alerted}{bg=light-bg}

  % ---------------------------------------------------------------------------
  % Formal title page: serif (Libertine, via \rmfamily) for title-page
  % elements only. Body text, itemize, and frametitle are intentionally left
  % on Beamer's sans-serif default -- \usebeamerfont{title} is also what
  % \transitionslide (below) uses for its big centered text, so section
  % transitions pick up the same serif treatment for free.
  % ---------------------------------------------------------------------------
  \setbeamerfont{title}{family=\rmfamily, series=\bfseries, size=\Large}
  \setbeamerfont{subtitle}{family=\rmfamily, shape=\itshape, size=\normalsize}
  \setbeamerfont{author}{family=\rmfamily, size=\normalsize}
  \setbeamerfont{institute}{family=\rmfamily, size=\scriptsize}
  \setbeamerfont{date}{family=\rmfamily, size=\scriptsize}

  \setbeamertemplate{title page}{
    \vfill
    \begin{center}
      {\usebeamerfont{title}\usebeamercolor[fg]{title}\inserttitle\par}
      \vspace{0.15cm}
      {\usebeamerfont{subtitle}\usebeamercolor[fg]{subtitle}\insertsubtitle\par}
      \vspace{0.45cm}
      {\color{primary-gold}\rule{0.42\textwidth}{0.6pt}\par}
      \vspace{0.45cm}
      {\usebeamerfont{author}\usebeamercolor[fg]{author}\insertauthor\par}
      \vspace{0.35cm}
      {\usebeamerfont{institute}\usebeamercolor[fg]{institute}\insertinstitute\par}
      \vspace{0.35cm}
      {\usebeamerfont{date}\usebeamercolor[fg]{date}\insertdate\par}
    \end{center}
    \vfill
  }

  % Minimal footer: course code (left) + page X / N (right), no navigation clutter
  \setbeamertemplate{navigation symbols}{}
  \setbeamertemplate{footline}{%
    \color{neutral}\footnotesize\hspace{0.7em}\@coursecodetext\hfill
    \insertframenumber/\inserttotalframenumber\hspace{0.7em}\vspace{0.3em}}
}{}
\makeatother

% -----------------------------------------------------------------------------
% TikZ setup — libraries the snippet gallery uses
% -----------------------------------------------------------------------------
\usepackage{tikz}
\usetikzlibrary{arrows.meta, positioning, calc, decorations.pathreplacing,
                fit, shapes.geometric, backgrounds}

% Shared TikZ styles — snippets and hand-written diagrams can pick these up
% via `\begin{tikzpicture}[every node/.style={...}]` or by naming specific
% styles (e.g., `\node[dag-node] (X) at (0,0) {X};`).
\tikzset{
  % Node styles for causal DAGs and similar
  dag-node/.style = {
    draw=primary-blue, fill=light-bg, rounded corners=2pt,
    minimum width=1.2cm, minimum height=0.9cm, align=center, font=\small
  },
  decision-node/.style = {
    draw=primary-gold, fill=white, diamond, aspect=1.8,
    minimum width=1.8cm, minimum height=1.2cm, align=center, font=\small,
    inner sep=1pt
  },
  % Edge styles — observed vs counterfactual
  observed-edge/.style = {-{Latex[length=2.5mm]}, thick, draw=jet},
  counterfactual-edge/.style = {-{Latex[length=2.5mm]}, thick, draw=neutral, dashed},
  confound-edge/.style = {-{Latex[length=2.5mm]}, thick, draw=negative, dashed},
  % Data-point styles for plots
  observed-dot/.style = {fill=primary-blue, draw=primary-blue, circle, inner sep=1.2pt},
  counterfactual-dot/.style = {fill=white, draw=neutral, circle, inner sep=1.2pt}
}

% -----------------------------------------------------------------------------
% Convenience macros
% -----------------------------------------------------------------------------
\newcommand{\muted}[1]{\textcolor{neutral}{#1}}
\newcommand{\key}[1]{\textcolor{primary-gold}{\textbf{#1}}}
\newcommand{\good}[1]{\textcolor{positive}{#1}}
\newcommand{\bad}[1]{\textcolor{negative}{#1}}

% Standout frame macro: full-bleed dark background for section transitions.
% Beamer-only (uses \begin{frame}/\setbeamercolor/\usebeamerfont) -- do not
% call from an article-class file (e.g. Notes/<CODE>/*-notes.tex). \muted,
% \key, \good, \bad, and \coursecode above are plain xcolor/gdef and are
% safe to use from either class; verified when Notes/article-class support
% was added.
% Expand: \transitionslide{Your transition title}
\newcommand{\transitionslide}[1]{%
  \begingroup
  \setbeamercolor{background canvas}{bg=primary-blue}
  \begin{frame}[plain]
    \centering
    \vfill
    {\usebeamerfont{title}\color{white}\Large #1\par}
    \vfill
  \end{frame}
  \endgroup
}

% and compiling with TEXINPUTS=../Preambles:$TEXINPUTS (the /compile-latex
% skill does this automatically).
%
% PALETTE CONTRACT. Color names defined here MUST match the names in
% Quarto/theme-template.scss so Beamer and Quarto renderings use the same
% palette. The scripts/check-palette-sync.sh script enforces this.
%
% To customize: edit the HEX values below AND the matching $variable in
% Quarto/theme-template.scss, then run scripts/check-palette-sync.sh.
%
% TITLE PAGE. A formal, serif (Libertine) title-page template is defined in
% the Beamer-only block below (\setbeamertemplate{title page}) -- title,
% subtitle, a thin gold rule, then \author/\institute/\date. Frametitle and
% body text remain on Beamer's sans-serif default. No SCSS counterpart is
% needed for this (font/layout only, no new colors).
% =============================================================================

% -----------------------------------------------------------------------------
% Engine detection + encoding
%
% Our /compile-latex skill uses XeLaTeX, where inputenc is unnecessary and
% can produce encoding warnings. Only load inputenc under pdfTeX.
% -----------------------------------------------------------------------------
\usepackage{iftex}
\ifPDFTeX
  \usepackage[utf8]{inputenc}
\fi

\usepackage{amsmath, amssymb, amsthm}
\usepackage{booktabs}
\usepackage{graphicx}

% -----------------------------------------------------------------------------
% Serif face for the title page (formal/academic look). Linux Libertine is
% XeLaTeX- and pdfTeX-safe and redefines \rmdefault, so \rmfamily picks it up
% everywhere \rmfamily is used explicitly. Body text and frametitles stay on
% Beamer's own sans-serif default (unaffected) -- only the title-page
% \setbeamerfont{...} overrides below opt in to \rmfamily.
% -----------------------------------------------------------------------------
\usepackage{libertine}

% -----------------------------------------------------------------------------
% PALETTE — mirrors Quarto/theme-template.scss variable names.
% Load xcolor and define colors BEFORE anything that references them
% (notably hyperref, hypersetup, and the Beamer color assignments).
% -----------------------------------------------------------------------------
\usepackage{xcolor}

\definecolor{primary-blue}{HTML}{012169}      % headings, accents
\definecolor{primary-gold}{HTML}{B9975B}      % emphasis, borders
\definecolor{highlight-yellow}{HTML}{F2A900}  % markers, alerts
\definecolor{light-bg}{HTML}{E8EDF5}          % title / section backgrounds
\definecolor{jet}{HTML}{1A1A1A}               % body text

% Semantic colors (used in diagrams and annotations)
\definecolor{positive}{HTML}{15803D}          % good / observed / identified
\definecolor{negative}{HTML}{B91C1C}          % bad / problematic / bias
\definecolor{neutral}{HTML}{525252}           % reference / context

% Highlight accents (match .hi-* classes in the SCSS)
\definecolor{hi-slate}{HTML}{314F4F}
\definecolor{hi-green}{HTML}{2E8B57}
\definecolor{hi-red}{HTML}{C41E3A}

% Semantic aliases so skill templates and snippets can write
% \textcolor{accent}{...} without hard-coding "primary-blue".
\colorlet{accent}{primary-blue}
\colorlet{accent2}{primary-gold}

% -----------------------------------------------------------------------------
% Hyperref — loaded AFTER xcolor + palette so link colors resolve.
% -----------------------------------------------------------------------------
\usepackage{hyperref}
\hypersetup{colorlinks=true, linkcolor=primary-blue, urlcolor=primary-blue,
            citecolor=primary-blue, pdfborder={0 0 0}}

% -----------------------------------------------------------------------------
% Course code — set once per deck (after % =============================================================================
% Preambles/header.tex — shared LaTeX/Beamer preamble for this project
%
% Use from a Beamer lecture by prepending:
%     \input{header}
% and compiling with TEXINPUTS=../Preambles:$TEXINPUTS (the /compile-latex
% skill does this automatically).
%
% PALETTE CONTRACT. Color names defined here MUST match the names in
% Quarto/theme-template.scss so Beamer and Quarto renderings use the same
% palette. The scripts/check-palette-sync.sh script enforces this.
%
% To customize: edit the HEX values below AND the matching $variable in
% Quarto/theme-template.scss, then run scripts/check-palette-sync.sh.
%
% TITLE PAGE. A formal, serif (Libertine) title-page template is defined in
% the Beamer-only block below (\setbeamertemplate{title page}) -- title,
% subtitle, a thin gold rule, then \author/\institute/\date. Frametitle and
% body text remain on Beamer's sans-serif default. No SCSS counterpart is
% needed for this (font/layout only, no new colors).
% =============================================================================

% -----------------------------------------------------------------------------
% Engine detection + encoding
%
% Our /compile-latex skill uses XeLaTeX, where inputenc is unnecessary and
% can produce encoding warnings. Only load inputenc under pdfTeX.
% -----------------------------------------------------------------------------
\usepackage{iftex}
\ifPDFTeX
  \usepackage[utf8]{inputenc}
\fi

\usepackage{amsmath, amssymb, amsthm}
\usepackage{booktabs}
\usepackage{graphicx}

% -----------------------------------------------------------------------------
% Serif face for the title page (formal/academic look). Linux Libertine is
% XeLaTeX- and pdfTeX-safe and redefines \rmdefault, so \rmfamily picks it up
% everywhere \rmfamily is used explicitly. Body text and frametitles stay on
% Beamer's own sans-serif default (unaffected) -- only the title-page
% \setbeamerfont{...} overrides below opt in to \rmfamily.
% -----------------------------------------------------------------------------
\usepackage{libertine}

% -----------------------------------------------------------------------------
% PALETTE — mirrors Quarto/theme-template.scss variable names.
% Load xcolor and define colors BEFORE anything that references them
% (notably hyperref, hypersetup, and the Beamer color assignments).
% -----------------------------------------------------------------------------
\usepackage{xcolor}

\definecolor{primary-blue}{HTML}{012169}      % headings, accents
\definecolor{primary-gold}{HTML}{B9975B}      % emphasis, borders
\definecolor{highlight-yellow}{HTML}{F2A900}  % markers, alerts
\definecolor{light-bg}{HTML}{E8EDF5}          % title / section backgrounds
\definecolor{jet}{HTML}{1A1A1A}               % body text

% Semantic colors (used in diagrams and annotations)
\definecolor{positive}{HTML}{15803D}          % good / observed / identified
\definecolor{negative}{HTML}{B91C1C}          % bad / problematic / bias
\definecolor{neutral}{HTML}{525252}           % reference / context

% Highlight accents (match .hi-* classes in the SCSS)
\definecolor{hi-slate}{HTML}{314F4F}
\definecolor{hi-green}{HTML}{2E8B57}
\definecolor{hi-red}{HTML}{C41E3A}

% Semantic aliases so skill templates and snippets can write
% \textcolor{accent}{...} without hard-coding "primary-blue".
\colorlet{accent}{primary-blue}
\colorlet{accent2}{primary-gold}

% -----------------------------------------------------------------------------
% Hyperref — loaded AFTER xcolor + palette so link colors resolve.
% -----------------------------------------------------------------------------
\usepackage{hyperref}
\hypersetup{colorlinks=true, linkcolor=primary-blue, urlcolor=primary-blue,
            citecolor=primary-blue, pdfborder={0 0 0}}

% -----------------------------------------------------------------------------
% Course code — set once per deck (after \input{header}, before \begin{document})
% via \coursecode{CS401}. Shown in the footer on every slide so a deck's course
% is identifiable without opening the title page. Empty by default.
% -----------------------------------------------------------------------------
\makeatletter
\newcommand{\coursecode}[1]{\gdef\@coursecodetext{#1}}
\gdef\@coursecodetext{}
\makeatother

% -----------------------------------------------------------------------------
% Beamer theme assignments (only apply under Beamer).
%
% We detect Beamer via \@ifclassloaded{beamer}, which is the documented,
% non-internal way. This must be wrapped in \makeatletter/\makeatother
% because \@ifclassloaded contains an @-catcode character.
% -----------------------------------------------------------------------------
\makeatletter
\@ifclassloaded{beamer}{%
  \usefonttheme{professionalfonts}
  \setbeamercolor{structure}{fg=primary-blue}
  \setbeamercolor{frametitle}{fg=primary-blue}
  \setbeamercolor{title}{fg=primary-blue}
  \setbeamercolor{subtitle}{fg=primary-gold}
  \setbeamercolor{author}{fg=primary-blue}
  \setbeamercolor{institute}{fg=neutral}
  \setbeamercolor{date}{fg=primary-gold}
  \setbeamercolor{itemize item}{fg=primary-blue}
  \setbeamercolor{itemize subitem}{fg=primary-gold}
  \setbeamercolor{alerted text}{fg=primary-gold}
  \setbeamercolor{block title}{fg=white, bg=primary-blue}
  \setbeamercolor{block body}{bg=light-bg}
  % Semantic block variants: block=definition/concept (blue), exampleblock=
  % worked example (green/positive), alertblock=key takeaway or caution (gold).
  % Keep to <=2 colored boxes per slide across ALL three variants (INV-7).
  \setbeamercolor{block title example}{fg=white, bg=positive}
  \setbeamercolor{block body example}{bg=light-bg}
  \setbeamercolor{block title alerted}{fg=white, bg=primary-gold}
  \setbeamercolor{block body alerted}{bg=light-bg}

  % ---------------------------------------------------------------------------
  % Formal title page: serif (Libertine, via \rmfamily) for title-page
  % elements only. Body text, itemize, and frametitle are intentionally left
  % on Beamer's sans-serif default -- \usebeamerfont{title} is also what
  % \transitionslide (below) uses for its big centered text, so section
  % transitions pick up the same serif treatment for free.
  % ---------------------------------------------------------------------------
  \setbeamerfont{title}{family=\rmfamily, series=\bfseries, size=\Large}
  \setbeamerfont{subtitle}{family=\rmfamily, shape=\itshape, size=\normalsize}
  \setbeamerfont{author}{family=\rmfamily, size=\normalsize}
  \setbeamerfont{institute}{family=\rmfamily, size=\scriptsize}
  \setbeamerfont{date}{family=\rmfamily, size=\scriptsize}

  \setbeamertemplate{title page}{
    \vfill
    \begin{center}
      {\usebeamerfont{title}\usebeamercolor[fg]{title}\inserttitle\par}
      \vspace{0.15cm}
      {\usebeamerfont{subtitle}\usebeamercolor[fg]{subtitle}\insertsubtitle\par}
      \vspace{0.45cm}
      {\color{primary-gold}\rule{0.42\textwidth}{0.6pt}\par}
      \vspace{0.45cm}
      {\usebeamerfont{author}\usebeamercolor[fg]{author}\insertauthor\par}
      \vspace{0.35cm}
      {\usebeamerfont{institute}\usebeamercolor[fg]{institute}\insertinstitute\par}
      \vspace{0.35cm}
      {\usebeamerfont{date}\usebeamercolor[fg]{date}\insertdate\par}
    \end{center}
    \vfill
  }

  % Minimal footer: course code (left) + page X / N (right), no navigation clutter
  \setbeamertemplate{navigation symbols}{}
  \setbeamertemplate{footline}{%
    \color{neutral}\footnotesize\hspace{0.7em}\@coursecodetext\hfill
    \insertframenumber/\inserttotalframenumber\hspace{0.7em}\vspace{0.3em}}
}{}
\makeatother

% -----------------------------------------------------------------------------
% TikZ setup — libraries the snippet gallery uses
% -----------------------------------------------------------------------------
\usepackage{tikz}
\usetikzlibrary{arrows.meta, positioning, calc, decorations.pathreplacing,
                fit, shapes.geometric, backgrounds}

% Shared TikZ styles — snippets and hand-written diagrams can pick these up
% via `\begin{tikzpicture}[every node/.style={...}]` or by naming specific
% styles (e.g., `\node[dag-node] (X) at (0,0) {X};`).
\tikzset{
  % Node styles for causal DAGs and similar
  dag-node/.style = {
    draw=primary-blue, fill=light-bg, rounded corners=2pt,
    minimum width=1.2cm, minimum height=0.9cm, align=center, font=\small
  },
  decision-node/.style = {
    draw=primary-gold, fill=white, diamond, aspect=1.8,
    minimum width=1.8cm, minimum height=1.2cm, align=center, font=\small,
    inner sep=1pt
  },
  % Edge styles — observed vs counterfactual
  observed-edge/.style = {-{Latex[length=2.5mm]}, thick, draw=jet},
  counterfactual-edge/.style = {-{Latex[length=2.5mm]}, thick, draw=neutral, dashed},
  confound-edge/.style = {-{Latex[length=2.5mm]}, thick, draw=negative, dashed},
  % Data-point styles for plots
  observed-dot/.style = {fill=primary-blue, draw=primary-blue, circle, inner sep=1.2pt},
  counterfactual-dot/.style = {fill=white, draw=neutral, circle, inner sep=1.2pt}
}

% -----------------------------------------------------------------------------
% Convenience macros
% -----------------------------------------------------------------------------
\newcommand{\muted}[1]{\textcolor{neutral}{#1}}
\newcommand{\key}[1]{\textcolor{primary-gold}{\textbf{#1}}}
\newcommand{\good}[1]{\textcolor{positive}{#1}}
\newcommand{\bad}[1]{\textcolor{negative}{#1}}

% Standout frame macro: full-bleed dark background for section transitions.
% Beamer-only (uses \begin{frame}/\setbeamercolor/\usebeamerfont) -- do not
% call from an article-class file (e.g. Notes/<CODE>/*-notes.tex). \muted,
% \key, \good, \bad, and \coursecode above are plain xcolor/gdef and are
% safe to use from either class; verified when Notes/article-class support
% was added.
% Expand: \transitionslide{Your transition title}
\newcommand{\transitionslide}[1]{%
  \begingroup
  \setbeamercolor{background canvas}{bg=primary-blue}
  \begin{frame}[plain]
    \centering
    \vfill
    {\usebeamerfont{title}\color{white}\Large #1\par}
    \vfill
  \end{frame}
  \endgroup
}
, before \begin{document})
% via \coursecode{CS401}. Shown in the footer on every slide so a deck's course
% is identifiable without opening the title page. Empty by default.
% -----------------------------------------------------------------------------
\makeatletter
\newcommand{\coursecode}[1]{\gdef\@coursecodetext{#1}}
\gdef\@coursecodetext{}
\makeatother

% -----------------------------------------------------------------------------
% Beamer theme assignments (only apply under Beamer).
%
% We detect Beamer via \@ifclassloaded{beamer}, which is the documented,
% non-internal way. This must be wrapped in \makeatletter/\makeatother
% because \@ifclassloaded contains an @-catcode character.
% -----------------------------------------------------------------------------
\makeatletter
\@ifclassloaded{beamer}{%
  \usefonttheme{professionalfonts}
  \setbeamercolor{structure}{fg=primary-blue}
  \setbeamercolor{frametitle}{fg=primary-blue}
  \setbeamercolor{title}{fg=primary-blue}
  \setbeamercolor{subtitle}{fg=primary-gold}
  \setbeamercolor{author}{fg=primary-blue}
  \setbeamercolor{institute}{fg=neutral}
  \setbeamercolor{date}{fg=primary-gold}
  \setbeamercolor{itemize item}{fg=primary-blue}
  \setbeamercolor{itemize subitem}{fg=primary-gold}
  \setbeamercolor{alerted text}{fg=primary-gold}
  \setbeamercolor{block title}{fg=white, bg=primary-blue}
  \setbeamercolor{block body}{bg=light-bg}
  % Semantic block variants: block=definition/concept (blue), exampleblock=
  % worked example (green/positive), alertblock=key takeaway or caution (gold).
  % Keep to <=2 colored boxes per slide across ALL three variants (INV-7).
  \setbeamercolor{block title example}{fg=white, bg=positive}
  \setbeamercolor{block body example}{bg=light-bg}
  \setbeamercolor{block title alerted}{fg=white, bg=primary-gold}
  \setbeamercolor{block body alerted}{bg=light-bg}

  % ---------------------------------------------------------------------------
  % Formal title page: serif (Libertine, via \rmfamily) for title-page
  % elements only. Body text, itemize, and frametitle are intentionally left
  % on Beamer's sans-serif default -- \usebeamerfont{title} is also what
  % \transitionslide (below) uses for its big centered text, so section
  % transitions pick up the same serif treatment for free.
  % ---------------------------------------------------------------------------
  \setbeamerfont{title}{family=\rmfamily, series=\bfseries, size=\Large}
  \setbeamerfont{subtitle}{family=\rmfamily, shape=\itshape, size=\normalsize}
  \setbeamerfont{author}{family=\rmfamily, size=\normalsize}
  \setbeamerfont{institute}{family=\rmfamily, size=\scriptsize}
  \setbeamerfont{date}{family=\rmfamily, size=\scriptsize}

  \setbeamertemplate{title page}{
    \vfill
    \begin{center}
      {\usebeamerfont{title}\usebeamercolor[fg]{title}\inserttitle\par}
      \vspace{0.15cm}
      {\usebeamerfont{subtitle}\usebeamercolor[fg]{subtitle}\insertsubtitle\par}
      \vspace{0.45cm}
      {\color{primary-gold}\rule{0.42\textwidth}{0.6pt}\par}
      \vspace{0.45cm}
      {\usebeamerfont{author}\usebeamercolor[fg]{author}\insertauthor\par}
      \vspace{0.35cm}
      {\usebeamerfont{institute}\usebeamercolor[fg]{institute}\insertinstitute\par}
      \vspace{0.35cm}
      {\usebeamerfont{date}\usebeamercolor[fg]{date}\insertdate\par}
    \end{center}
    \vfill
  }

  % Minimal footer: course code (left) + page X / N (right), no navigation clutter
  \setbeamertemplate{navigation symbols}{}
  \setbeamertemplate{footline}{%
    \color{neutral}\footnotesize\hspace{0.7em}\@coursecodetext\hfill
    \insertframenumber/\inserttotalframenumber\hspace{0.7em}\vspace{0.3em}}
}{}
\makeatother

% -----------------------------------------------------------------------------
% TikZ setup — libraries the snippet gallery uses
% -----------------------------------------------------------------------------
\usepackage{tikz}
\usetikzlibrary{arrows.meta, positioning, calc, decorations.pathreplacing,
                fit, shapes.geometric, backgrounds}

% Shared TikZ styles — snippets and hand-written diagrams can pick these up
% via `\begin{tikzpicture}[every node/.style={...}]` or by naming specific
% styles (e.g., `\node[dag-node] (X) at (0,0) {X};`).
\tikzset{
  % Node styles for causal DAGs and similar
  dag-node/.style = {
    draw=primary-blue, fill=light-bg, rounded corners=2pt,
    minimum width=1.2cm, minimum height=0.9cm, align=center, font=\small
  },
  decision-node/.style = {
    draw=primary-gold, fill=white, diamond, aspect=1.8,
    minimum width=1.8cm, minimum height=1.2cm, align=center, font=\small,
    inner sep=1pt
  },
  % Edge styles — observed vs counterfactual
  observed-edge/.style = {-{Latex[length=2.5mm]}, thick, draw=jet},
  counterfactual-edge/.style = {-{Latex[length=2.5mm]}, thick, draw=neutral, dashed},
  confound-edge/.style = {-{Latex[length=2.5mm]}, thick, draw=negative, dashed},
  % Data-point styles for plots
  observed-dot/.style = {fill=primary-blue, draw=primary-blue, circle, inner sep=1.2pt},
  counterfactual-dot/.style = {fill=white, draw=neutral, circle, inner sep=1.2pt}
}

% -----------------------------------------------------------------------------
% Convenience macros
% -----------------------------------------------------------------------------
\newcommand{\muted}[1]{\textcolor{neutral}{#1}}
\newcommand{\key}[1]{\textcolor{primary-gold}{\textbf{#1}}}
\newcommand{\good}[1]{\textcolor{positive}{#1}}
\newcommand{\bad}[1]{\textcolor{negative}{#1}}

% Standout frame macro: full-bleed dark background for section transitions.
% Beamer-only (uses \begin{frame}/\setbeamercolor/\usebeamerfont) -- do not
% call from an article-class file (e.g. Notes/<CODE>/*-notes.tex). \muted,
% \key, \good, \bad, and \coursecode above are plain xcolor/gdef and are
% safe to use from either class; verified when Notes/article-class support
% was added.
% Expand: \transitionslide{Your transition title}
\newcommand{\transitionslide}[1]{%
  \begingroup
  \setbeamercolor{background canvas}{bg=primary-blue}
  \begin{frame}[plain]
    \centering
    \vfill
    {\usebeamerfont{title}\color{white}\Large #1\par}
    \vfill
  \end{frame}
  \endgroup
}

% and compiling with TEXINPUTS=../Preambles:$TEXINPUTS (the /compile-latex
% skill does this automatically).
%
% PALETTE CONTRACT. Color names defined here MUST match the names in
% Quarto/theme-template.scss so Beamer and Quarto renderings use the same
% palette. The scripts/check-palette-sync.sh script enforces this.
%
% To customize: edit the HEX values below AND the matching $variable in
% Quarto/theme-template.scss, then run scripts/check-palette-sync.sh.
%
% TITLE PAGE. A formal, serif (Libertine) title-page template is defined in
% the Beamer-only block below (\setbeamertemplate{title page}) -- title,
% subtitle, a thin gold rule, then \author/\institute/\date. Frametitle and
% body text remain on Beamer's sans-serif default. No SCSS counterpart is
% needed for this (font/layout only, no new colors).
% =============================================================================

% -----------------------------------------------------------------------------
% Engine detection + encoding
%
% Our /compile-latex skill uses XeLaTeX, where inputenc is unnecessary and
% can produce encoding warnings. Only load inputenc under pdfTeX.
% -----------------------------------------------------------------------------
\usepackage{iftex}
\ifPDFTeX
  \usepackage[utf8]{inputenc}
\fi

\usepackage{amsmath, amssymb, amsthm}
\usepackage{booktabs}
\usepackage{graphicx}

% -----------------------------------------------------------------------------
% Serif face for the title page (formal/academic look). Linux Libertine is
% XeLaTeX- and pdfTeX-safe and redefines \rmdefault, so \rmfamily picks it up
% everywhere \rmfamily is used explicitly. Body text and frametitles stay on
% Beamer's own sans-serif default (unaffected) -- only the title-page
% \setbeamerfont{...} overrides below opt in to \rmfamily.
% -----------------------------------------------------------------------------
\usepackage{libertine}

% -----------------------------------------------------------------------------
% PALETTE — mirrors Quarto/theme-template.scss variable names.
% Load xcolor and define colors BEFORE anything that references them
% (notably hyperref, hypersetup, and the Beamer color assignments).
% -----------------------------------------------------------------------------
\usepackage{xcolor}

\definecolor{primary-blue}{HTML}{012169}      % headings, accents
\definecolor{primary-gold}{HTML}{B9975B}      % emphasis, borders
\definecolor{highlight-yellow}{HTML}{F2A900}  % markers, alerts
\definecolor{light-bg}{HTML}{E8EDF5}          % title / section backgrounds
\definecolor{jet}{HTML}{1A1A1A}               % body text

% Semantic colors (used in diagrams and annotations)
\definecolor{positive}{HTML}{15803D}          % good / observed / identified
\definecolor{negative}{HTML}{B91C1C}          % bad / problematic / bias
\definecolor{neutral}{HTML}{525252}           % reference / context

% Highlight accents (match .hi-* classes in the SCSS)
\definecolor{hi-slate}{HTML}{314F4F}
\definecolor{hi-green}{HTML}{2E8B57}
\definecolor{hi-red}{HTML}{C41E3A}

% Semantic aliases so skill templates and snippets can write
% \textcolor{accent}{...} without hard-coding "primary-blue".
\colorlet{accent}{primary-blue}
\colorlet{accent2}{primary-gold}

% -----------------------------------------------------------------------------
% Hyperref — loaded AFTER xcolor + palette so link colors resolve.
% -----------------------------------------------------------------------------
\usepackage{hyperref}
\hypersetup{colorlinks=true, linkcolor=primary-blue, urlcolor=primary-blue,
            citecolor=primary-blue, pdfborder={0 0 0}}

% -----------------------------------------------------------------------------
% Course code — set once per deck (after % =============================================================================
% Preambles/header.tex — shared LaTeX/Beamer preamble for this project
%
% Use from a Beamer lecture by prepending:
%     % =============================================================================
% Preambles/header.tex — shared LaTeX/Beamer preamble for this project
%
% Use from a Beamer lecture by prepending:
%     \input{header}
% and compiling with TEXINPUTS=../Preambles:$TEXINPUTS (the /compile-latex
% skill does this automatically).
%
% PALETTE CONTRACT. Color names defined here MUST match the names in
% Quarto/theme-template.scss so Beamer and Quarto renderings use the same
% palette. The scripts/check-palette-sync.sh script enforces this.
%
% To customize: edit the HEX values below AND the matching $variable in
% Quarto/theme-template.scss, then run scripts/check-palette-sync.sh.
%
% TITLE PAGE. A formal, serif (Libertine) title-page template is defined in
% the Beamer-only block below (\setbeamertemplate{title page}) -- title,
% subtitle, a thin gold rule, then \author/\institute/\date. Frametitle and
% body text remain on Beamer's sans-serif default. No SCSS counterpart is
% needed for this (font/layout only, no new colors).
% =============================================================================

% -----------------------------------------------------------------------------
% Engine detection + encoding
%
% Our /compile-latex skill uses XeLaTeX, where inputenc is unnecessary and
% can produce encoding warnings. Only load inputenc under pdfTeX.
% -----------------------------------------------------------------------------
\usepackage{iftex}
\ifPDFTeX
  \usepackage[utf8]{inputenc}
\fi

\usepackage{amsmath, amssymb, amsthm}
\usepackage{booktabs}
\usepackage{graphicx}

% -----------------------------------------------------------------------------
% Serif face for the title page (formal/academic look). Linux Libertine is
% XeLaTeX- and pdfTeX-safe and redefines \rmdefault, so \rmfamily picks it up
% everywhere \rmfamily is used explicitly. Body text and frametitles stay on
% Beamer's own sans-serif default (unaffected) -- only the title-page
% \setbeamerfont{...} overrides below opt in to \rmfamily.
% -----------------------------------------------------------------------------
\usepackage{libertine}

% -----------------------------------------------------------------------------
% PALETTE — mirrors Quarto/theme-template.scss variable names.
% Load xcolor and define colors BEFORE anything that references them
% (notably hyperref, hypersetup, and the Beamer color assignments).
% -----------------------------------------------------------------------------
\usepackage{xcolor}

\definecolor{primary-blue}{HTML}{012169}      % headings, accents
\definecolor{primary-gold}{HTML}{B9975B}      % emphasis, borders
\definecolor{highlight-yellow}{HTML}{F2A900}  % markers, alerts
\definecolor{light-bg}{HTML}{E8EDF5}          % title / section backgrounds
\definecolor{jet}{HTML}{1A1A1A}               % body text

% Semantic colors (used in diagrams and annotations)
\definecolor{positive}{HTML}{15803D}          % good / observed / identified
\definecolor{negative}{HTML}{B91C1C}          % bad / problematic / bias
\definecolor{neutral}{HTML}{525252}           % reference / context

% Highlight accents (match .hi-* classes in the SCSS)
\definecolor{hi-slate}{HTML}{314F4F}
\definecolor{hi-green}{HTML}{2E8B57}
\definecolor{hi-red}{HTML}{C41E3A}

% Semantic aliases so skill templates and snippets can write
% \textcolor{accent}{...} without hard-coding "primary-blue".
\colorlet{accent}{primary-blue}
\colorlet{accent2}{primary-gold}

% -----------------------------------------------------------------------------
% Hyperref — loaded AFTER xcolor + palette so link colors resolve.
% -----------------------------------------------------------------------------
\usepackage{hyperref}
\hypersetup{colorlinks=true, linkcolor=primary-blue, urlcolor=primary-blue,
            citecolor=primary-blue, pdfborder={0 0 0}}

% -----------------------------------------------------------------------------
% Course code — set once per deck (after \input{header}, before \begin{document})
% via \coursecode{CS401}. Shown in the footer on every slide so a deck's course
% is identifiable without opening the title page. Empty by default.
% -----------------------------------------------------------------------------
\makeatletter
\newcommand{\coursecode}[1]{\gdef\@coursecodetext{#1}}
\gdef\@coursecodetext{}
\makeatother

% -----------------------------------------------------------------------------
% Beamer theme assignments (only apply under Beamer).
%
% We detect Beamer via \@ifclassloaded{beamer}, which is the documented,
% non-internal way. This must be wrapped in \makeatletter/\makeatother
% because \@ifclassloaded contains an @-catcode character.
% -----------------------------------------------------------------------------
\makeatletter
\@ifclassloaded{beamer}{%
  \usefonttheme{professionalfonts}
  \setbeamercolor{structure}{fg=primary-blue}
  \setbeamercolor{frametitle}{fg=primary-blue}
  \setbeamercolor{title}{fg=primary-blue}
  \setbeamercolor{subtitle}{fg=primary-gold}
  \setbeamercolor{author}{fg=primary-blue}
  \setbeamercolor{institute}{fg=neutral}
  \setbeamercolor{date}{fg=primary-gold}
  \setbeamercolor{itemize item}{fg=primary-blue}
  \setbeamercolor{itemize subitem}{fg=primary-gold}
  \setbeamercolor{alerted text}{fg=primary-gold}
  \setbeamercolor{block title}{fg=white, bg=primary-blue}
  \setbeamercolor{block body}{bg=light-bg}
  % Semantic block variants: block=definition/concept (blue), exampleblock=
  % worked example (green/positive), alertblock=key takeaway or caution (gold).
  % Keep to <=2 colored boxes per slide across ALL three variants (INV-7).
  \setbeamercolor{block title example}{fg=white, bg=positive}
  \setbeamercolor{block body example}{bg=light-bg}
  \setbeamercolor{block title alerted}{fg=white, bg=primary-gold}
  \setbeamercolor{block body alerted}{bg=light-bg}

  % ---------------------------------------------------------------------------
  % Formal title page: serif (Libertine, via \rmfamily) for title-page
  % elements only. Body text, itemize, and frametitle are intentionally left
  % on Beamer's sans-serif default -- \usebeamerfont{title} is also what
  % \transitionslide (below) uses for its big centered text, so section
  % transitions pick up the same serif treatment for free.
  % ---------------------------------------------------------------------------
  \setbeamerfont{title}{family=\rmfamily, series=\bfseries, size=\Large}
  \setbeamerfont{subtitle}{family=\rmfamily, shape=\itshape, size=\normalsize}
  \setbeamerfont{author}{family=\rmfamily, size=\normalsize}
  \setbeamerfont{institute}{family=\rmfamily, size=\scriptsize}
  \setbeamerfont{date}{family=\rmfamily, size=\scriptsize}

  \setbeamertemplate{title page}{
    \vfill
    \begin{center}
      {\usebeamerfont{title}\usebeamercolor[fg]{title}\inserttitle\par}
      \vspace{0.15cm}
      {\usebeamerfont{subtitle}\usebeamercolor[fg]{subtitle}\insertsubtitle\par}
      \vspace{0.45cm}
      {\color{primary-gold}\rule{0.42\textwidth}{0.6pt}\par}
      \vspace{0.45cm}
      {\usebeamerfont{author}\usebeamercolor[fg]{author}\insertauthor\par}
      \vspace{0.35cm}
      {\usebeamerfont{institute}\usebeamercolor[fg]{institute}\insertinstitute\par}
      \vspace{0.35cm}
      {\usebeamerfont{date}\usebeamercolor[fg]{date}\insertdate\par}
    \end{center}
    \vfill
  }

  % Minimal footer: course code (left) + page X / N (right), no navigation clutter
  \setbeamertemplate{navigation symbols}{}
  \setbeamertemplate{footline}{%
    \color{neutral}\footnotesize\hspace{0.7em}\@coursecodetext\hfill
    \insertframenumber/\inserttotalframenumber\hspace{0.7em}\vspace{0.3em}}
}{}
\makeatother

% -----------------------------------------------------------------------------
% TikZ setup — libraries the snippet gallery uses
% -----------------------------------------------------------------------------
\usepackage{tikz}
\usetikzlibrary{arrows.meta, positioning, calc, decorations.pathreplacing,
                fit, shapes.geometric, backgrounds}

% Shared TikZ styles — snippets and hand-written diagrams can pick these up
% via `\begin{tikzpicture}[every node/.style={...}]` or by naming specific
% styles (e.g., `\node[dag-node] (X) at (0,0) {X};`).
\tikzset{
  % Node styles for causal DAGs and similar
  dag-node/.style = {
    draw=primary-blue, fill=light-bg, rounded corners=2pt,
    minimum width=1.2cm, minimum height=0.9cm, align=center, font=\small
  },
  decision-node/.style = {
    draw=primary-gold, fill=white, diamond, aspect=1.8,
    minimum width=1.8cm, minimum height=1.2cm, align=center, font=\small,
    inner sep=1pt
  },
  % Edge styles — observed vs counterfactual
  observed-edge/.style = {-{Latex[length=2.5mm]}, thick, draw=jet},
  counterfactual-edge/.style = {-{Latex[length=2.5mm]}, thick, draw=neutral, dashed},
  confound-edge/.style = {-{Latex[length=2.5mm]}, thick, draw=negative, dashed},
  % Data-point styles for plots
  observed-dot/.style = {fill=primary-blue, draw=primary-blue, circle, inner sep=1.2pt},
  counterfactual-dot/.style = {fill=white, draw=neutral, circle, inner sep=1.2pt}
}

% -----------------------------------------------------------------------------
% Convenience macros
% -----------------------------------------------------------------------------
\newcommand{\muted}[1]{\textcolor{neutral}{#1}}
\newcommand{\key}[1]{\textcolor{primary-gold}{\textbf{#1}}}
\newcommand{\good}[1]{\textcolor{positive}{#1}}
\newcommand{\bad}[1]{\textcolor{negative}{#1}}

% Standout frame macro: full-bleed dark background for section transitions.
% Beamer-only (uses \begin{frame}/\setbeamercolor/\usebeamerfont) -- do not
% call from an article-class file (e.g. Notes/<CODE>/*-notes.tex). \muted,
% \key, \good, \bad, and \coursecode above are plain xcolor/gdef and are
% safe to use from either class; verified when Notes/article-class support
% was added.
% Expand: \transitionslide{Your transition title}
\newcommand{\transitionslide}[1]{%
  \begingroup
  \setbeamercolor{background canvas}{bg=primary-blue}
  \begin{frame}[plain]
    \centering
    \vfill
    {\usebeamerfont{title}\color{white}\Large #1\par}
    \vfill
  \end{frame}
  \endgroup
}

% and compiling with TEXINPUTS=../Preambles:$TEXINPUTS (the /compile-latex
% skill does this automatically).
%
% PALETTE CONTRACT. Color names defined here MUST match the names in
% Quarto/theme-template.scss so Beamer and Quarto renderings use the same
% palette. The scripts/check-palette-sync.sh script enforces this.
%
% To customize: edit the HEX values below AND the matching $variable in
% Quarto/theme-template.scss, then run scripts/check-palette-sync.sh.
%
% TITLE PAGE. A formal, serif (Libertine) title-page template is defined in
% the Beamer-only block below (\setbeamertemplate{title page}) -- title,
% subtitle, a thin gold rule, then \author/\institute/\date. Frametitle and
% body text remain on Beamer's sans-serif default. No SCSS counterpart is
% needed for this (font/layout only, no new colors).
% =============================================================================

% -----------------------------------------------------------------------------
% Engine detection + encoding
%
% Our /compile-latex skill uses XeLaTeX, where inputenc is unnecessary and
% can produce encoding warnings. Only load inputenc under pdfTeX.
% -----------------------------------------------------------------------------
\usepackage{iftex}
\ifPDFTeX
  \usepackage[utf8]{inputenc}
\fi

\usepackage{amsmath, amssymb, amsthm}
\usepackage{booktabs}
\usepackage{graphicx}

% -----------------------------------------------------------------------------
% Serif face for the title page (formal/academic look). Linux Libertine is
% XeLaTeX- and pdfTeX-safe and redefines \rmdefault, so \rmfamily picks it up
% everywhere \rmfamily is used explicitly. Body text and frametitles stay on
% Beamer's own sans-serif default (unaffected) -- only the title-page
% \setbeamerfont{...} overrides below opt in to \rmfamily.
% -----------------------------------------------------------------------------
\usepackage{libertine}

% -----------------------------------------------------------------------------
% PALETTE — mirrors Quarto/theme-template.scss variable names.
% Load xcolor and define colors BEFORE anything that references them
% (notably hyperref, hypersetup, and the Beamer color assignments).
% -----------------------------------------------------------------------------
\usepackage{xcolor}

\definecolor{primary-blue}{HTML}{012169}      % headings, accents
\definecolor{primary-gold}{HTML}{B9975B}      % emphasis, borders
\definecolor{highlight-yellow}{HTML}{F2A900}  % markers, alerts
\definecolor{light-bg}{HTML}{E8EDF5}          % title / section backgrounds
\definecolor{jet}{HTML}{1A1A1A}               % body text

% Semantic colors (used in diagrams and annotations)
\definecolor{positive}{HTML}{15803D}          % good / observed / identified
\definecolor{negative}{HTML}{B91C1C}          % bad / problematic / bias
\definecolor{neutral}{HTML}{525252}           % reference / context

% Highlight accents (match .hi-* classes in the SCSS)
\definecolor{hi-slate}{HTML}{314F4F}
\definecolor{hi-green}{HTML}{2E8B57}
\definecolor{hi-red}{HTML}{C41E3A}

% Semantic aliases so skill templates and snippets can write
% \textcolor{accent}{...} without hard-coding "primary-blue".
\colorlet{accent}{primary-blue}
\colorlet{accent2}{primary-gold}

% -----------------------------------------------------------------------------
% Hyperref — loaded AFTER xcolor + palette so link colors resolve.
% -----------------------------------------------------------------------------
\usepackage{hyperref}
\hypersetup{colorlinks=true, linkcolor=primary-blue, urlcolor=primary-blue,
            citecolor=primary-blue, pdfborder={0 0 0}}

% -----------------------------------------------------------------------------
% Course code — set once per deck (after % =============================================================================
% Preambles/header.tex — shared LaTeX/Beamer preamble for this project
%
% Use from a Beamer lecture by prepending:
%     \input{header}
% and compiling with TEXINPUTS=../Preambles:$TEXINPUTS (the /compile-latex
% skill does this automatically).
%
% PALETTE CONTRACT. Color names defined here MUST match the names in
% Quarto/theme-template.scss so Beamer and Quarto renderings use the same
% palette. The scripts/check-palette-sync.sh script enforces this.
%
% To customize: edit the HEX values below AND the matching $variable in
% Quarto/theme-template.scss, then run scripts/check-palette-sync.sh.
%
% TITLE PAGE. A formal, serif (Libertine) title-page template is defined in
% the Beamer-only block below (\setbeamertemplate{title page}) -- title,
% subtitle, a thin gold rule, then \author/\institute/\date. Frametitle and
% body text remain on Beamer's sans-serif default. No SCSS counterpart is
% needed for this (font/layout only, no new colors).
% =============================================================================

% -----------------------------------------------------------------------------
% Engine detection + encoding
%
% Our /compile-latex skill uses XeLaTeX, where inputenc is unnecessary and
% can produce encoding warnings. Only load inputenc under pdfTeX.
% -----------------------------------------------------------------------------
\usepackage{iftex}
\ifPDFTeX
  \usepackage[utf8]{inputenc}
\fi

\usepackage{amsmath, amssymb, amsthm}
\usepackage{booktabs}
\usepackage{graphicx}

% -----------------------------------------------------------------------------
% Serif face for the title page (formal/academic look). Linux Libertine is
% XeLaTeX- and pdfTeX-safe and redefines \rmdefault, so \rmfamily picks it up
% everywhere \rmfamily is used explicitly. Body text and frametitles stay on
% Beamer's own sans-serif default (unaffected) -- only the title-page
% \setbeamerfont{...} overrides below opt in to \rmfamily.
% -----------------------------------------------------------------------------
\usepackage{libertine}

% -----------------------------------------------------------------------------
% PALETTE — mirrors Quarto/theme-template.scss variable names.
% Load xcolor and define colors BEFORE anything that references them
% (notably hyperref, hypersetup, and the Beamer color assignments).
% -----------------------------------------------------------------------------
\usepackage{xcolor}

\definecolor{primary-blue}{HTML}{012169}      % headings, accents
\definecolor{primary-gold}{HTML}{B9975B}      % emphasis, borders
\definecolor{highlight-yellow}{HTML}{F2A900}  % markers, alerts
\definecolor{light-bg}{HTML}{E8EDF5}          % title / section backgrounds
\definecolor{jet}{HTML}{1A1A1A}               % body text

% Semantic colors (used in diagrams and annotations)
\definecolor{positive}{HTML}{15803D}          % good / observed / identified
\definecolor{negative}{HTML}{B91C1C}          % bad / problematic / bias
\definecolor{neutral}{HTML}{525252}           % reference / context

% Highlight accents (match .hi-* classes in the SCSS)
\definecolor{hi-slate}{HTML}{314F4F}
\definecolor{hi-green}{HTML}{2E8B57}
\definecolor{hi-red}{HTML}{C41E3A}

% Semantic aliases so skill templates and snippets can write
% \textcolor{accent}{...} without hard-coding "primary-blue".
\colorlet{accent}{primary-blue}
\colorlet{accent2}{primary-gold}

% -----------------------------------------------------------------------------
% Hyperref — loaded AFTER xcolor + palette so link colors resolve.
% -----------------------------------------------------------------------------
\usepackage{hyperref}
\hypersetup{colorlinks=true, linkcolor=primary-blue, urlcolor=primary-blue,
            citecolor=primary-blue, pdfborder={0 0 0}}

% -----------------------------------------------------------------------------
% Course code — set once per deck (after \input{header}, before \begin{document})
% via \coursecode{CS401}. Shown in the footer on every slide so a deck's course
% is identifiable without opening the title page. Empty by default.
% -----------------------------------------------------------------------------
\makeatletter
\newcommand{\coursecode}[1]{\gdef\@coursecodetext{#1}}
\gdef\@coursecodetext{}
\makeatother

% -----------------------------------------------------------------------------
% Beamer theme assignments (only apply under Beamer).
%
% We detect Beamer via \@ifclassloaded{beamer}, which is the documented,
% non-internal way. This must be wrapped in \makeatletter/\makeatother
% because \@ifclassloaded contains an @-catcode character.
% -----------------------------------------------------------------------------
\makeatletter
\@ifclassloaded{beamer}{%
  \usefonttheme{professionalfonts}
  \setbeamercolor{structure}{fg=primary-blue}
  \setbeamercolor{frametitle}{fg=primary-blue}
  \setbeamercolor{title}{fg=primary-blue}
  \setbeamercolor{subtitle}{fg=primary-gold}
  \setbeamercolor{author}{fg=primary-blue}
  \setbeamercolor{institute}{fg=neutral}
  \setbeamercolor{date}{fg=primary-gold}
  \setbeamercolor{itemize item}{fg=primary-blue}
  \setbeamercolor{itemize subitem}{fg=primary-gold}
  \setbeamercolor{alerted text}{fg=primary-gold}
  \setbeamercolor{block title}{fg=white, bg=primary-blue}
  \setbeamercolor{block body}{bg=light-bg}
  % Semantic block variants: block=definition/concept (blue), exampleblock=
  % worked example (green/positive), alertblock=key takeaway or caution (gold).
  % Keep to <=2 colored boxes per slide across ALL three variants (INV-7).
  \setbeamercolor{block title example}{fg=white, bg=positive}
  \setbeamercolor{block body example}{bg=light-bg}
  \setbeamercolor{block title alerted}{fg=white, bg=primary-gold}
  \setbeamercolor{block body alerted}{bg=light-bg}

  % ---------------------------------------------------------------------------
  % Formal title page: serif (Libertine, via \rmfamily) for title-page
  % elements only. Body text, itemize, and frametitle are intentionally left
  % on Beamer's sans-serif default -- \usebeamerfont{title} is also what
  % \transitionslide (below) uses for its big centered text, so section
  % transitions pick up the same serif treatment for free.
  % ---------------------------------------------------------------------------
  \setbeamerfont{title}{family=\rmfamily, series=\bfseries, size=\Large}
  \setbeamerfont{subtitle}{family=\rmfamily, shape=\itshape, size=\normalsize}
  \setbeamerfont{author}{family=\rmfamily, size=\normalsize}
  \setbeamerfont{institute}{family=\rmfamily, size=\scriptsize}
  \setbeamerfont{date}{family=\rmfamily, size=\scriptsize}

  \setbeamertemplate{title page}{
    \vfill
    \begin{center}
      {\usebeamerfont{title}\usebeamercolor[fg]{title}\inserttitle\par}
      \vspace{0.15cm}
      {\usebeamerfont{subtitle}\usebeamercolor[fg]{subtitle}\insertsubtitle\par}
      \vspace{0.45cm}
      {\color{primary-gold}\rule{0.42\textwidth}{0.6pt}\par}
      \vspace{0.45cm}
      {\usebeamerfont{author}\usebeamercolor[fg]{author}\insertauthor\par}
      \vspace{0.35cm}
      {\usebeamerfont{institute}\usebeamercolor[fg]{institute}\insertinstitute\par}
      \vspace{0.35cm}
      {\usebeamerfont{date}\usebeamercolor[fg]{date}\insertdate\par}
    \end{center}
    \vfill
  }

  % Minimal footer: course code (left) + page X / N (right), no navigation clutter
  \setbeamertemplate{navigation symbols}{}
  \setbeamertemplate{footline}{%
    \color{neutral}\footnotesize\hspace{0.7em}\@coursecodetext\hfill
    \insertframenumber/\inserttotalframenumber\hspace{0.7em}\vspace{0.3em}}
}{}
\makeatother

% -----------------------------------------------------------------------------
% TikZ setup — libraries the snippet gallery uses
% -----------------------------------------------------------------------------
\usepackage{tikz}
\usetikzlibrary{arrows.meta, positioning, calc, decorations.pathreplacing,
                fit, shapes.geometric, backgrounds}

% Shared TikZ styles — snippets and hand-written diagrams can pick these up
% via `\begin{tikzpicture}[every node/.style={...}]` or by naming specific
% styles (e.g., `\node[dag-node] (X) at (0,0) {X};`).
\tikzset{
  % Node styles for causal DAGs and similar
  dag-node/.style = {
    draw=primary-blue, fill=light-bg, rounded corners=2pt,
    minimum width=1.2cm, minimum height=0.9cm, align=center, font=\small
  },
  decision-node/.style = {
    draw=primary-gold, fill=white, diamond, aspect=1.8,
    minimum width=1.8cm, minimum height=1.2cm, align=center, font=\small,
    inner sep=1pt
  },
  % Edge styles — observed vs counterfactual
  observed-edge/.style = {-{Latex[length=2.5mm]}, thick, draw=jet},
  counterfactual-edge/.style = {-{Latex[length=2.5mm]}, thick, draw=neutral, dashed},
  confound-edge/.style = {-{Latex[length=2.5mm]}, thick, draw=negative, dashed},
  % Data-point styles for plots
  observed-dot/.style = {fill=primary-blue, draw=primary-blue, circle, inner sep=1.2pt},
  counterfactual-dot/.style = {fill=white, draw=neutral, circle, inner sep=1.2pt}
}

% -----------------------------------------------------------------------------
% Convenience macros
% -----------------------------------------------------------------------------
\newcommand{\muted}[1]{\textcolor{neutral}{#1}}
\newcommand{\key}[1]{\textcolor{primary-gold}{\textbf{#1}}}
\newcommand{\good}[1]{\textcolor{positive}{#1}}
\newcommand{\bad}[1]{\textcolor{negative}{#1}}

% Standout frame macro: full-bleed dark background for section transitions.
% Beamer-only (uses \begin{frame}/\setbeamercolor/\usebeamerfont) -- do not
% call from an article-class file (e.g. Notes/<CODE>/*-notes.tex). \muted,
% \key, \good, \bad, and \coursecode above are plain xcolor/gdef and are
% safe to use from either class; verified when Notes/article-class support
% was added.
% Expand: \transitionslide{Your transition title}
\newcommand{\transitionslide}[1]{%
  \begingroup
  \setbeamercolor{background canvas}{bg=primary-blue}
  \begin{frame}[plain]
    \centering
    \vfill
    {\usebeamerfont{title}\color{white}\Large #1\par}
    \vfill
  \end{frame}
  \endgroup
}
, before \begin{document})
% via \coursecode{CS401}. Shown in the footer on every slide so a deck's course
% is identifiable without opening the title page. Empty by default.
% -----------------------------------------------------------------------------
\makeatletter
\newcommand{\coursecode}[1]{\gdef\@coursecodetext{#1}}
\gdef\@coursecodetext{}
\makeatother

% -----------------------------------------------------------------------------
% Beamer theme assignments (only apply under Beamer).
%
% We detect Beamer via \@ifclassloaded{beamer}, which is the documented,
% non-internal way. This must be wrapped in \makeatletter/\makeatother
% because \@ifclassloaded contains an @-catcode character.
% -----------------------------------------------------------------------------
\makeatletter
\@ifclassloaded{beamer}{%
  \usefonttheme{professionalfonts}
  \setbeamercolor{structure}{fg=primary-blue}
  \setbeamercolor{frametitle}{fg=primary-blue}
  \setbeamercolor{title}{fg=primary-blue}
  \setbeamercolor{subtitle}{fg=primary-gold}
  \setbeamercolor{author}{fg=primary-blue}
  \setbeamercolor{institute}{fg=neutral}
  \setbeamercolor{date}{fg=primary-gold}
  \setbeamercolor{itemize item}{fg=primary-blue}
  \setbeamercolor{itemize subitem}{fg=primary-gold}
  \setbeamercolor{alerted text}{fg=primary-gold}
  \setbeamercolor{block title}{fg=white, bg=primary-blue}
  \setbeamercolor{block body}{bg=light-bg}
  % Semantic block variants: block=definition/concept (blue), exampleblock=
  % worked example (green/positive), alertblock=key takeaway or caution (gold).
  % Keep to <=2 colored boxes per slide across ALL three variants (INV-7).
  \setbeamercolor{block title example}{fg=white, bg=positive}
  \setbeamercolor{block body example}{bg=light-bg}
  \setbeamercolor{block title alerted}{fg=white, bg=primary-gold}
  \setbeamercolor{block body alerted}{bg=light-bg}

  % ---------------------------------------------------------------------------
  % Formal title page: serif (Libertine, via \rmfamily) for title-page
  % elements only. Body text, itemize, and frametitle are intentionally left
  % on Beamer's sans-serif default -- \usebeamerfont{title} is also what
  % \transitionslide (below) uses for its big centered text, so section
  % transitions pick up the same serif treatment for free.
  % ---------------------------------------------------------------------------
  \setbeamerfont{title}{family=\rmfamily, series=\bfseries, size=\Large}
  \setbeamerfont{subtitle}{family=\rmfamily, shape=\itshape, size=\normalsize}
  \setbeamerfont{author}{family=\rmfamily, size=\normalsize}
  \setbeamerfont{institute}{family=\rmfamily, size=\scriptsize}
  \setbeamerfont{date}{family=\rmfamily, size=\scriptsize}

  \setbeamertemplate{title page}{
    \vfill
    \begin{center}
      {\usebeamerfont{title}\usebeamercolor[fg]{title}\inserttitle\par}
      \vspace{0.15cm}
      {\usebeamerfont{subtitle}\usebeamercolor[fg]{subtitle}\insertsubtitle\par}
      \vspace{0.45cm}
      {\color{primary-gold}\rule{0.42\textwidth}{0.6pt}\par}
      \vspace{0.45cm}
      {\usebeamerfont{author}\usebeamercolor[fg]{author}\insertauthor\par}
      \vspace{0.35cm}
      {\usebeamerfont{institute}\usebeamercolor[fg]{institute}\insertinstitute\par}
      \vspace{0.35cm}
      {\usebeamerfont{date}\usebeamercolor[fg]{date}\insertdate\par}
    \end{center}
    \vfill
  }

  % Minimal footer: course code (left) + page X / N (right), no navigation clutter
  \setbeamertemplate{navigation symbols}{}
  \setbeamertemplate{footline}{%
    \color{neutral}\footnotesize\hspace{0.7em}\@coursecodetext\hfill
    \insertframenumber/\inserttotalframenumber\hspace{0.7em}\vspace{0.3em}}
}{}
\makeatother

% -----------------------------------------------------------------------------
% TikZ setup — libraries the snippet gallery uses
% -----------------------------------------------------------------------------
\usepackage{tikz}
\usetikzlibrary{arrows.meta, positioning, calc, decorations.pathreplacing,
                fit, shapes.geometric, backgrounds}

% Shared TikZ styles — snippets and hand-written diagrams can pick these up
% via `\begin{tikzpicture}[every node/.style={...}]` or by naming specific
% styles (e.g., `\node[dag-node] (X) at (0,0) {X};`).
\tikzset{
  % Node styles for causal DAGs and similar
  dag-node/.style = {
    draw=primary-blue, fill=light-bg, rounded corners=2pt,
    minimum width=1.2cm, minimum height=0.9cm, align=center, font=\small
  },
  decision-node/.style = {
    draw=primary-gold, fill=white, diamond, aspect=1.8,
    minimum width=1.8cm, minimum height=1.2cm, align=center, font=\small,
    inner sep=1pt
  },
  % Edge styles — observed vs counterfactual
  observed-edge/.style = {-{Latex[length=2.5mm]}, thick, draw=jet},
  counterfactual-edge/.style = {-{Latex[length=2.5mm]}, thick, draw=neutral, dashed},
  confound-edge/.style = {-{Latex[length=2.5mm]}, thick, draw=negative, dashed},
  % Data-point styles for plots
  observed-dot/.style = {fill=primary-blue, draw=primary-blue, circle, inner sep=1.2pt},
  counterfactual-dot/.style = {fill=white, draw=neutral, circle, inner sep=1.2pt}
}

% -----------------------------------------------------------------------------
% Convenience macros
% -----------------------------------------------------------------------------
\newcommand{\muted}[1]{\textcolor{neutral}{#1}}
\newcommand{\key}[1]{\textcolor{primary-gold}{\textbf{#1}}}
\newcommand{\good}[1]{\textcolor{positive}{#1}}
\newcommand{\bad}[1]{\textcolor{negative}{#1}}

% Standout frame macro: full-bleed dark background for section transitions.
% Beamer-only (uses \begin{frame}/\setbeamercolor/\usebeamerfont) -- do not
% call from an article-class file (e.g. Notes/<CODE>/*-notes.tex). \muted,
% \key, \good, \bad, and \coursecode above are plain xcolor/gdef and are
% safe to use from either class; verified when Notes/article-class support
% was added.
% Expand: \transitionslide{Your transition title}
\newcommand{\transitionslide}[1]{%
  \begingroup
  \setbeamercolor{background canvas}{bg=primary-blue}
  \begin{frame}[plain]
    \centering
    \vfill
    {\usebeamerfont{title}\color{white}\Large #1\par}
    \vfill
  \end{frame}
  \endgroup
}
, before \begin{document})
% via \coursecode{CS401}. Shown in the footer on every slide so a deck's course
% is identifiable without opening the title page. Empty by default.
% -----------------------------------------------------------------------------
\makeatletter
\newcommand{\coursecode}[1]{\gdef\@coursecodetext{#1}}
\gdef\@coursecodetext{}
\makeatother

% -----------------------------------------------------------------------------
% Beamer theme assignments (only apply under Beamer).
%
% We detect Beamer via \@ifclassloaded{beamer}, which is the documented,
% non-internal way. This must be wrapped in \makeatletter/\makeatother
% because \@ifclassloaded contains an @-catcode character.
% -----------------------------------------------------------------------------
\makeatletter
\@ifclassloaded{beamer}{%
  \usefonttheme{professionalfonts}
  \setbeamercolor{structure}{fg=primary-blue}
  \setbeamercolor{frametitle}{fg=primary-blue}
  \setbeamercolor{title}{fg=primary-blue}
  \setbeamercolor{subtitle}{fg=primary-gold}
  \setbeamercolor{author}{fg=primary-blue}
  \setbeamercolor{institute}{fg=neutral}
  \setbeamercolor{date}{fg=primary-gold}
  \setbeamercolor{itemize item}{fg=primary-blue}
  \setbeamercolor{itemize subitem}{fg=primary-gold}
  \setbeamercolor{alerted text}{fg=primary-gold}
  \setbeamercolor{block title}{fg=white, bg=primary-blue}
  \setbeamercolor{block body}{bg=light-bg}
  % Semantic block variants: block=definition/concept (blue), exampleblock=
  % worked example (green/positive), alertblock=key takeaway or caution (gold).
  % Keep to <=2 colored boxes per slide across ALL three variants (INV-7).
  \setbeamercolor{block title example}{fg=white, bg=positive}
  \setbeamercolor{block body example}{bg=light-bg}
  \setbeamercolor{block title alerted}{fg=white, bg=primary-gold}
  \setbeamercolor{block body alerted}{bg=light-bg}

  % ---------------------------------------------------------------------------
  % Formal title page: serif (Libertine, via \rmfamily) for title-page
  % elements only. Body text, itemize, and frametitle are intentionally left
  % on Beamer's sans-serif default -- \usebeamerfont{title} is also what
  % \transitionslide (below) uses for its big centered text, so section
  % transitions pick up the same serif treatment for free.
  % ---------------------------------------------------------------------------
  \setbeamerfont{title}{family=\rmfamily, series=\bfseries, size=\Large}
  \setbeamerfont{subtitle}{family=\rmfamily, shape=\itshape, size=\normalsize}
  \setbeamerfont{author}{family=\rmfamily, size=\normalsize}
  \setbeamerfont{institute}{family=\rmfamily, size=\scriptsize}
  \setbeamerfont{date}{family=\rmfamily, size=\scriptsize}

  \setbeamertemplate{title page}{
    \vfill
    \begin{center}
      {\usebeamerfont{title}\usebeamercolor[fg]{title}\inserttitle\par}
      \vspace{0.15cm}
      {\usebeamerfont{subtitle}\usebeamercolor[fg]{subtitle}\insertsubtitle\par}
      \vspace{0.45cm}
      {\color{primary-gold}\rule{0.42\textwidth}{0.6pt}\par}
      \vspace{0.45cm}
      {\usebeamerfont{author}\usebeamercolor[fg]{author}\insertauthor\par}
      \vspace{0.35cm}
      {\usebeamerfont{institute}\usebeamercolor[fg]{institute}\insertinstitute\par}
      \vspace{0.35cm}
      {\usebeamerfont{date}\usebeamercolor[fg]{date}\insertdate\par}
    \end{center}
    \vfill
  }

  % Minimal footer: course code (left) + page X / N (right), no navigation clutter
  \setbeamertemplate{navigation symbols}{}
  \setbeamertemplate{footline}{%
    \color{neutral}\footnotesize\hspace{0.7em}\@coursecodetext\hfill
    \insertframenumber/\inserttotalframenumber\hspace{0.7em}\vspace{0.3em}}
}{}
\makeatother

% -----------------------------------------------------------------------------
% TikZ setup — libraries the snippet gallery uses
% -----------------------------------------------------------------------------
\usepackage{tikz}
\usetikzlibrary{arrows.meta, positioning, calc, decorations.pathreplacing,
                fit, shapes.geometric, backgrounds}

% Shared TikZ styles — snippets and hand-written diagrams can pick these up
% via `\begin{tikzpicture}[every node/.style={...}]` or by naming specific
% styles (e.g., `\node[dag-node] (X) at (0,0) {X};`).
\tikzset{
  % Node styles for causal DAGs and similar
  dag-node/.style = {
    draw=primary-blue, fill=light-bg, rounded corners=2pt,
    minimum width=1.2cm, minimum height=0.9cm, align=center, font=\small
  },
  decision-node/.style = {
    draw=primary-gold, fill=white, diamond, aspect=1.8,
    minimum width=1.8cm, minimum height=1.2cm, align=center, font=\small,
    inner sep=1pt
  },
  % Edge styles — observed vs counterfactual
  observed-edge/.style = {-{Latex[length=2.5mm]}, thick, draw=jet},
  counterfactual-edge/.style = {-{Latex[length=2.5mm]}, thick, draw=neutral, dashed},
  confound-edge/.style = {-{Latex[length=2.5mm]}, thick, draw=negative, dashed},
  % Data-point styles for plots
  observed-dot/.style = {fill=primary-blue, draw=primary-blue, circle, inner sep=1.2pt},
  counterfactual-dot/.style = {fill=white, draw=neutral, circle, inner sep=1.2pt}
}

% -----------------------------------------------------------------------------
% Convenience macros
% -----------------------------------------------------------------------------
\newcommand{\muted}[1]{\textcolor{neutral}{#1}}
\newcommand{\key}[1]{\textcolor{primary-gold}{\textbf{#1}}}
\newcommand{\good}[1]{\textcolor{positive}{#1}}
\newcommand{\bad}[1]{\textcolor{negative}{#1}}

% Standout frame macro: full-bleed dark background for section transitions.
% Beamer-only (uses \begin{frame}/\setbeamercolor/\usebeamerfont) -- do not
% call from an article-class file (e.g. Notes/<CODE>/*-notes.tex). \muted,
% \key, \good, \bad, and \coursecode above are plain xcolor/gdef and are
% safe to use from either class; verified when Notes/article-class support
% was added.
% Expand: \transitionslide{Your transition title}
\newcommand{\transitionslide}[1]{%
  \begingroup
  \setbeamercolor{background canvas}{bg=primary-blue}
  \begin{frame}[plain]
    \centering
    \vfill
    {\usebeamerfont{title}\color{white}\Large #1\par}
    \vfill
  \end{frame}
  \endgroup
}

\coursecode{CS401}
\renewcommand{\thesection}{8.\arabic{section}}

\title{Competitive Exam Practice 8 --- Answer Key\\\large I/O Techniques}
\author{Auqib Hamid Lone\\[0.3em]
  \emph{Department of Computer Science \& Engineering}, \emph{GCET Ganderbal}}
\date{\today}

\begin{document}
\maketitle

\begin{enumerate}
\item[\textbf{Q1}] \textbf{(D).} The CPU never reacts mid-instruction: the lecture's five-step scenario (step 2) states it finishes its \emph{current instruction} first, then checks \texttt{PS.IE}, then branches to the ISR. (A) and (B) describe behaviours no interrupt should cause; (C) omits the finish-the-instruction condition.

\item[\textbf{Q2}] \textbf{(B).} In a vectored interrupt each device supplies its own identifying code, and hardware uses it to index the \textbf{interrupt-vector table} --- a table of ISR starting addresses, one entry per device --- jumping straight to the correct ISR. (A) describes a fixed vector location (not device-supplied), (C) describes polling a processor register, and (D) is wrong because (B) holds.

\item[\textbf{Q3}] \textbf{(A).} The lecture's two-switch gating makes this precise: a request reaches the CPU only when the device's enable \emph{and} \texttt{PS.IE} both allow it, so with interrupts disabled the CPU will not process them. (C) is false as stated in GATE's intended sense --- the CPU checks for interrupts \emph{after} completing an instruction, i.e.\ between instructions, not ``before executing a new instruction'' in the way (C) claims to be the distinguishing truth; (B) and (D) are plainly false (interruptible loops exist; edge-triggered interrupts exist).

\item[\textbf{Q4}] \textbf{(A).} In a daisy chain, the bus-grant (or interrupt-acknowledge) signal propagates device-to-device in a fixed physical order, so the device closest to the processor is granted first --- priority is non-uniform (highest for the nearest), not uniform (B), and it needs only one shared signal path, not a separate pin per device (D).

\item[\textbf{Q5}] \textbf{TRUE.} A bus (DMA) request outranks an interrupt request: the DMA/bus request is the lower-latency, hardware-initiated path that must be served before the interrupt can even be processed (the interrupt itself needs the bus), and it cannot be delayed without risking loss of the on-going transfer. This matches the lecture's ordering of techniques, where DMA sits \emph{below} the interrupt layer on the critical path.

\item[\textbf{Q6}] \textbf{TRUE.} In memory-mapped I/O the device registers live in the normal address space, so \emph{any} memory-referencing instruction (including block-move instructions) can be used to transfer data; I/O-mapped I/O is restricted to dedicated \texttt{IN}/\texttt{OUT} instructions. More flexible, direct addressing --- no special instruction decode, no narrower address range --- makes data transfer between microprocessor and device \emph{usually} faster, which is the sense the statement asserts.

\item[\textbf{Q7}] \textbf{TRUE.} Programmed (polling) I/O, when the device is already ready, performs the data transfer directly --- a single read/write of the data register. Interrupt-driven I/O always carries the overhead of the full five-step handling scenario on top of the same transfer: finish the current instruction, save \texttt{PC} (and shadow registers), disable \texttt{IE}, jump to the ISR, service the device, and execute \texttt{RTI} to restore state and re-enable interrupts. That per-transfer overhead is why programmed I/O is faster \emph{per transfer} when the device is ready, exactly the lecture's poll-vs-interrupt trade-off table: polling's cost is the wasted waiting, not the transfer itself.

\item[\textbf{Q8}] \textbf{100000000} ($10^8$). One keystroke every $0.1$~s; at $2\times10^9$ instructions/s the CPU executes $0.1\times 2\times10^9=2\times10^8$ instructions between keystrokes; at 2 instructions per failed iteration, that is $2\times10^8 \div 2 = 10^8$ iterations --- each one re-reading an unchanged status flag, accomplishing nothing useful.

\item[\textbf{Q9}] \textbf{(B).} The synchronous/async trade-off table is decisive here: asynchronous pays \emph{two full round-trip delays} per transfer (steps 1--2 and 3--4 of the handshake), which a synchronous bus never pays for a consistently fast device; synchronous only loses on slow devices (wait states) or when the clock period must shrink to beat skew. For devices of known, similar speed, the cheaper synchronous protocol wins.

\item[\textbf{Q10}] \textbf{(B).} Cycle stealing is the DMA controller temporarily taking memory bus cycles the CPU would otherwise have used, because both compete for the same shared memory bus --- the CPU is delayed only for the individual stolen cycles, not halted. (C) inverts the whole point of DMA (the CPU is \emph{not} running one ISR per word), and (A)/(D) get the direction of the competition backwards.

\item[\textbf{Q11}] \textbf{(B).} Where DMA offloads \emph{moving} data (one pre-configured block transfer, then stop), an IOP is a dedicated processor given a whole \emph{channel program} --- several transfers, data formats, and simple decisions --- that it executes independently, interrupting the main CPU only when the entire program completes. (A), (C), and (D) contradict the lecture: an IOP is \emph{more} hardware than a DMA controller, does interrupt the CPU (once, at program completion), and must still arbitrate for the bus when it moves data.
\end{enumerate}

\end{document}
